\chapter{Vectori și valori proprii. Diagonalizare}

\section{Vectori și valori proprii}

Fie $ V $ un $ K $-spațiu vectorial și $ f : V \to V $ o aplicație
liniară.
\begin{definition}\label{def:vvp} \index{vector!propriu}
  Un vector $ v \in V $ se numește \emph{vector propriu} (eng.\ \emph{eigenvector})
  pentru aplicația $ f $ dacă există un scalar $ \lambda \in K $ astfel încît
  $ f(v) = \lambda v $.

  În acest caz, $ \lambda $ se numește \emph{valoarea proprie} (eng.\ \emph{eigenvalue})
  asociată vectorului propriu $ v $.
\end{definition}

Rezultă că vectorii proprii sînt aceia pentru care aplicația liniară
are o acțiune simplă, de \qq{rescalare}, adică doar de înmulțire cu un scalar,
care se numește valoarea proprie asociată.

Pașii pentru calculul vectorilor și valorilor proprii sînt:
\begin{enumerate}[(1)]
\item Presupunem $ \dim V = n $. Scriem matricea aplicației $ f $ în baza
  canonică a lui $ V $ și obținem $ A = M_f^B \in M_n(K) $;
\item Scriem \emph{polinomul caracteristic} al matricei, anume:
  \[
    P(x) = \det(A - x \cdot I_n).
  \]
\item Rădăcinile polinomului caracteristic sînt valorile proprii ale
  endomorfismului $ f $. Mulțimea valorilor proprii se mai numește \emph{spectrul}
  endomorfismului și se notează $ \sigma(f) $. \index{polinom!caracteristic}
  \index{valoare proprie!spectru}
\item Pentru a găsi vectorii proprii asociați fiecărei valori proprii $ \lambda_i $,
  se rezolvă ecuația $ f(v_i) = \lambda_i v_i $ și se determină $ v_i \in V $.
\end{enumerate}

Alte elemente de terminologie:
\begin{definition}\label{def:mult}
  Fie $ \lambda $ o valoare proprie a unui endomorfism $ f $, iar $ v $, vectorul
  propriu asociat.

  Notăm $ V(\lambda) = V_\lambda = \dr{Sp}(v) $ subspațiul lui $ V $ generat
  de $ v $, numit \emph{subspațiul invariant} (propriu) asociat lui $ v $.

  Dimensiunea acestui subspațiu se numește \emph{multiplicitatea geometrică}
  a valorii proprii, notată $ m_g(\lambda) = g(\lambda) = \dim V(\lambda) $.
  \index{valoare proprie!multiplicitate geometrică}
  \index{valoare proprie!multiplicitate algebrică}
  \index{vector!propriu!subspațiu invariant}

  Se numește \emph{multiplicitatea algebrică} a valorii proprii $ \lambda $,
  notată $ m_a(\lambda) = a(\lambda) $, multiplicitatea rădăcinii $ x = \lambda $
  în polinomul caracteristic. Adică
  $ a(\lambda) = n \Leftrightarrow P_A(x) \vdots (x - \lambda)^n $.
\end{definition}

O proprietate importantă este:
\begin{theorem}[Cayley-Hamilton]\label{thm:ch}
  \index{teorema!Cayley-Hamilton}
  $ P_A(A) = 0 $, unde $ A $ este polinomul caracteristic al matricei $ A $.
\end{theorem}

%%%%%%%%%%%%%%%%%%%%%%%%%%%%%%%%%%%%%%%%%%%%%%%%%%%%%%%%%%%%%%%%%%%%%%

\section{Matrice de trecere. Diagonalizare}
\index{matrice!de trecere}
\index{matrice!diagonalizare}

Putem avea două sau mai multe baze ale aceluiași spațiu vectorial,
iar între ele există o legătură strînsă, matriceală.

Fie $ B = \{ b_1, \dots, b_n \} $ și $ C = \{ c_1, \dots, c_n \} $ două
baza ale aceluiași spațiu vectorial $ V $. Se numește \emph{matricea de trecere}
de la baza $ B $ la baza $ C $, notată $ M_B^C $ sau $ {}^B M^C $, matricea
coeficienților din scrierea vectorilor $ c_i $ în funcție de vectorii $ b_i $.
Mai precis, deoarece $ B $ este bază, avem:
\[
  \forall i,j, \quad c_i = \sum_j \alpha_{ij} b_j,
\]
iar matricea de trecere este matricea $ (\alpha_{ij}) $.

\begin{definition}\label{def:diag}
  O matrice $ A \in M_n(K) $ se numește \emph{diagonalizabilă} dacă ea
  este asemenea cu o matrice diagonală $ D $, care are elemente nenule doar
  pe diagonala principală.

  Altfel spus, există $ T \in M_n(K) $ inversabilă, cu $ T^{-1} AT = D $.
\end{definition}

\begin{remark}\label{rk:diag}
  Dacă matricea $ A $ este diagonalizabilă, atunci $ D $ din definiția de mai sus
  conține pe diagonală doar valorile proprii ale lui $ A $.
\end{remark}

Condițiile necesare și suficiente pentru diagonalizare sînt:
\begin{theorem}\label{thm:diag}
  Următoarele afirmații sînt echivalente:
  \begin{enumerate}[(a)]
  \item Matricea $ A $ este diagonalizabilă;
  \item Există o bază a spațiului vectorial $ V $ formată doar din vectorii
    proprii ai lui $ A $;
  \item $ a(\lambda_i) = g(\lambda_i), \forall \lambda_i \in \sigma(A) $.
  \end{enumerate}
\end{theorem}

Evident, discuția are sens atît pentru cazul în care pornim cu o aplicație
liniară, caz în care îi luăm matricea în baza canonică și lucrăm cu ea,
cît și dacă pornim direct cu o matrice.

Pașii pentru diagonalizare sînt:
\begin{enumerate}[(1)]
\item Fixăm o bază $ B $ a lui $ V $ (e.g.\ baza canonică) și scriem $ A = M_f^B $;
\item Determinăm valorile proprii $ \lambda_i $ ale lui $ A $ și multiplicitățile
  lor geometrice;
\item Pentru fiecare valoare proprie, determinăm vectorii proprii, subspațiile
  invariante, cîte o bază $ B_i $ în fiecare dintre acestea și multiplicitățile
  geometrice;
\item Dacă există o valoare proprie $ \lambda_j $ cu $ a(\lambda_j) \neq g(\lambda_j)$,
  algoritmul se oprește și matricea nu se poate diagonaliza;
\item Dacă $ a(\lambda_i) = g(\lambda_i), \forall i $, matricea se poate
  diagonaliza, iar $ B' = \ds\cup_i B_i $ este o bază a lui $ V $, formată
  numai din vectori proprii;
\item Se determină matricea de trecere de la baza $ B $ la baza $ B' $, notată $ T $,
  care este inversabilă, iar forma diagonală este:
  \begin{equation}
    \label{eq:diag}
    A \sim T^{-1} AT = \dr{diag} (\lambda_i).
  \end{equation}
\end{enumerate}

Observați că ecuația~\eqref{eq:diag} poate fi scrisă și astfel:
\[
  T^{-1}AT = \dr{diag}(\lambda_i) \stackrel{\text{not.}}{=\joinrel=} D \Leftrightarrow %
  A = TDT^{-1},
\]
unde $ D $ este o matrice diagonală. Astfel că, în general, se poate formula următoarea:

\begin{definition}\label{def:diag}
  O matrice $ A \in M_n(\RR) $ este diagonalizabilă dacă există o matrice diagonală $ D $
  și o matrice inversabilă $ T $ astfel încît $ A = TDT^{-1} $.
\end{definition}

Această definiție abstractă este concretizată în procedura de mai sus:
\begin{itemize}
\item matricea inversabilă este matricea de trecere de la baza canonică la o bază
  formată doar din vectori proprii;
\item matricea diagonală conține valorile proprii pe diagonala principală.
\end{itemize}

\begin{remark}\label{rk:diag-putere}
  Unul dintre avantajele diagonalizării matricelor este ușurința calculelor
  ulterioare. În particular, dacă $ A \sim \dr{diag}(\lambda_i) $,
  atunci $ A^k \sim \dr{diag}(\lambda_i^k), \forall k $.

  Mai mult, folosind relația din ecuația~\eqref{eq:diag}, putem scrie și:
  \[
    A^n = (TDT^{-1})^n = TD^nT^{-1},
  \]
  cu avantajul că $ D^n = \left(\mathrm{diag}(\lambda_i)\right)^n = \mathrm{diag}(\lambda_i^n) $.
\end{remark}

%%%%%%%%%%%%%%%%%%%%%%%%%%%%%%%%%%%%%%%%%%%%%%%%%%%%%%%%%%%%%%%%%%%%%%

\section{Exerciții}

\indent\indent 1. Să se determine vectorii și valorile proprii pentru matricele:
\[
  A = \begin{pmatrix} 1 & 0 \\ 8 & 2 \end{pmatrix}, \quad %
  B = \begin{pmatrix} 0 & 2 & 2 \\ 2 & 0 & -2 \\ 2 & -2 & 0 \end{pmatrix}, \quad %
  C = \begin{pmatrix} 0 & 1 & -2 \\ 1 & 1 & 1 \\ 1 & -1 & -1 \end{pmatrix}.
\]

2. Fie aplicația liniară $ f : \RR^3 \to \RR^3, f(x, y, z) = (-y, x, z) $.

Să se calculeze vectorii și valorile proprii ale lui $ f $.

\vspace{0.3cm}

3. Aceeași cerință pentru:
\[
  f : \RR^3 \to \RR^3, \quad f(x, y, z) = (x - z, 8x + y - 2z, z).
\]

4. Fie aplicația liniară $ f : M_2(\RR) \to M_2(\RR) $, definită prin:
\[
  f \begin{pmatrix} a & b \\ c & d \end{pmatrix} = %
  \begin{pmatrix} 2a + b + d & 2b + 4c + 5d \\ 2c + d & 8d \end{pmatrix}.
\]

Să se arate că $ f $ nu este diagonalizabilă.

\vspace{0.3cm}

5. Să se diagonalizeze endomorfismul:
\[
  f : \RR^3 \to \RR^3, \quad f(x, y, z) = (y + z, x + z, x + y).
\]

6. Să se diagonalizeze matricele:
\[
  A = \begin{pmatrix}
    7 & 2 & -2 \\
    2 & 4 & -1 \\
    -2 & -1 & 4
  \end{pmatrix}, \quad %
  B = \begin{pmatrix}
    1 & 1 \\ 4 & 1
  \end{pmatrix}, \quad %
  C = \begin{pmatrix}
    11 & -5 & 5 \\
    -5 & 3 & -3 \\
    5 & -3 & 3
  \end{pmatrix}.
\]


7. Determinați vectorii și valorile proprii pentru aplicația:
\[
  f : \RR^3 \to \RR^3, \quad f(a, b, c) = (a - b + c, b, b-c).
\]

%%% Local Variables:
%%% mode: latex
%%% TeX-master: "../ag-20-21"
%%% End:
